\documentclass[10pt,a4paper]{amsart}
\usepackage{amsmath,amssymb,amsthm,graphicx}
% Retained minimal preamble for the Stacks Project chapter more-algebra.tex (More on Algebra).
% The chapter file pinned-original-source.tex carries only \input{preamble} on its first line, so no
% preamble text lives inside the chapter itself.  The Stacks Project uses one single shared preamble for
% every chapter; the snapshot stages it identically as sources/stacks-chow/preamble.tex and
% sources/stacks-etale-cohomology/preamble.tex.  The block below is copied mechanically, line for line,
% from sources/stacks-chow/preamble.tex lines 14-18 (the Xy-pic setup) and 154-217 (the native theorem
% declarations together with the mathematical macros).  Nothing was added, removed or redefined; the
% wrapper class and the amsmath/amssymb/amsthm/graphicx base packages are supplied by the assembler.
% Chapter: more-algebra.tex -- More on Algebra.
% For commutative diagrams we use Xy-pic
\usepackage[all]{xy}

% We use 2cell for 2-commutative diagrams.
\xyoption{2cell}

% Theorem environments.
%
\theoremstyle{plain}
\newtheorem{theorem}[subsection]{Theorem}
\newtheorem{proposition}[subsection]{Proposition}
\newtheorem{lemma}[subsection]{Lemma}

\theoremstyle{definition}
\newtheorem{definition}[subsection]{Definition}
\newtheorem{example}[subsection]{Example}
\newtheorem{exercise}[subsection]{Exercise}
\newtheorem{situation}[subsection]{Situation}

\theoremstyle{remark}
\newtheorem{remark}[subsection]{Remark}
\newtheorem{remarks}[subsection]{Remarks}

\numberwithin{equation}{subsection}

% Macros
%
\def\lim{\mathop{\mathrm{lim}}\nolimits}
\def\colim{\mathop{\mathrm{colim}}\nolimits}
\def\Spec{\mathop{\mathrm{Spec}}}
\def\Hom{\mathop{\mathrm{Hom}}\nolimits}
\def\Ext{\mathop{\mathrm{Ext}}\nolimits}
\def\SheafHom{\mathop{\mathcal{H}\!\mathit{om}}\nolimits}
\def\SheafExt{\mathop{\mathcal{E}\!\mathit{xt}}\nolimits}
\def\Sch{\mathit{Sch}}
\def\Mor{\mathop{\mathrm{Mor}}\nolimits}
\def\Ob{\mathop{\mathrm{Ob}}\nolimits}
\def\Sh{\mathop{\mathit{Sh}}\nolimits}
\def\NL{\mathop{N\!L}\nolimits}
\def\CH{\mathop{\mathrm{CH}}\nolimits}
\def\proetale{{pro\text{-}\acute{e}tale}}
\def\etale{{\acute{e}tale}}
\def\QCoh{\mathit{QCoh}}
\def\Ker{\mathop{\mathrm{Ker}}}
\def\Im{\mathop{\mathrm{Im}}}
\def\Coker{\mathop{\mathrm{Coker}}}
\def\Coim{\mathop{\mathrm{Coim}}}

% Boxtimes
%
\DeclareMathSymbol{\boxtimes}{\mathbin}{AMSa}{"02}

%
% Macros for moduli stacks/spaces
%
\def\QCohstack{\mathcal{QC}\!\mathit{oh}}
\def\Cohstack{\mathcal{C}\!\mathit{oh}}
\def\Spacesstack{\mathcal{S}\!\mathit{paces}}
\def\Quotfunctor{\mathrm{Quot}}
\def\Hilbfunctor{\mathrm{Hilb}}
\def\Curvesstack{\mathcal{C}\!\mathit{urves}}
\def\Polarizedstack{\mathcal{P}\!\mathit{olarized}}
\def\Complexesstack{\mathcal{C}\!\mathit{omplexes}}
% \Pic is the operator that assigns to X its picard group, usage \Pic(X)
% \Picardstack_{X/B} denotes the Picard stack of X over B
% \Picardfunctor_{X/B} denotes the Picard functor of X over B
\def\Pic{\mathop{\mathrm{Pic}}\nolimits}
\def\Picardstack{\mathcal{P}\!\mathit{ic}}
\def\Picardfunctor{\mathrm{Pic}}
\def\Deformationcategory{\mathcal{D}\!\mathit{ef}}
\usepackage{cleveref}
\crefname{originallemma}{original lemma}{original lemmas}
\Crefname{originallemma}{Original lemma}{Original lemmas}
\makeatletter
\newlabel{more-algebra-lemma-finite-module-over-fibre-product}{{0D2H}{source}{}{sourcepointer-1}{}}
\newlabel{more-algebra-lemma-finite-module-over-fibre-product@cref}{{[originallemma][0][]0D2H}{[1][0][]source}{}{}{}}
\makeatother
\crefname{originallemma}{original lemma}{original lemmas}
\Crefname{originallemma}{Original lemma}{Original lemmas}
\makeatletter
\newlabel{lemma-module-over-fibre-product}{{07RU}{source}{}{sourcepointer-2}{}}
\newlabel{lemma-module-over-fibre-product@cref}{{[originallemma][0][]07RU}{[1][0][]source}{}{}{}}
\makeatother
\title{Finiteness of the fibre product of two finite modules over a fibre product of rings}
\author{The Stacks Project Authors}
\date{stacks-project:more-algebra.tex:local-snapshot-20261012; GNU FDL 1.2 or later}
\begin{document}
\maketitle
\setcounter{subsection}{1}
\noindent\textit{Editorial note.} All delivered mathematical text is the author's own, copied from the pinned Stacks Project chapter file more-algebra.tex; no statement or proof step was rewritten, reordered or removed.\par\medskip
\noindent\textit{Editorial note.} The chapter file pinned-original-source.tex carries only \textbackslash{}input\{preamble\} on its first line, so no preamble text lives inside the chapter. The retained file sources/stacks-more-algebra/retained-preamble.tex is a mechanical, line-for-line copy of the project's single shared preamble as staged in the snapshot (Xy-pic setup plus theorem declarations and mathematical macros). Nothing was added, removed or redefined, and the excerpt serves only as preamble, never as a span coordinate.\par\medskip
\noindent\textit{Editorial note.} Both counted spans start at the \textbackslash{}begin line of their environment, so the statement and the author's proof are delivered whole; surrounding prose of the section is deliberately outside the counted spans and is carried as standard prerequisites instead.\par\medskip
\noindent\textit{Editorial note.} The citation fields (source\_id, source\_version, authors, source\_url, license.name, source\_title) are identical to those of every other card this lane has produced from this chapter, so the batch selector treats them as one source.\par\medskip
\noindent\textit{Allowed standard prerequisites (editorial).}\begin{itemize}
\item The Situation of a fibre product of rings and the identification of the two base changes of the fibre product module with the outer modules. the author cites Lemma lemma-module-over-fibre-product, counted by the sibling card S1-algebra-module-over-fibre-product at 643-715, at proof lines 774-776. Supplies the identifications used to lift the generators.
\item Surjectivity of the structure map from the fibre product of rings onto each factor and the identification of its kernel with the given ideal. used at proof lines 782-790 and 798-802. Lets coefficients in the ideal be lifted to elements of the fibre product.
\item The fibre product of modules consists of pairs agreeing in the common quotient. used at proof lines 794-802 to recognise the difference of two elements as a pair vanishing in the first component.. Supplies the data and identifications used in the corresponding step of the proof.
\end{itemize}\medskip
\ifdefined\hypertarget\hypertarget{sourcepointer-1}{}\fi
\noindent\textit{Source reference.} Original lemma 0D2H: focused-seeds/AGP/sources/stacks-snapshot-20261011/tags.txt, source lines 16382--16382. The lemma's own canonical tag, quoted so that readers can resolve the label inside the delivered text; it points at the counted statement itself and carries no proof content.\par\medskip
\ifdefined\hypertarget\hypertarget{sourcepointer-2}{}\fi
\noindent\textit{Source reference.} Original lemma 07RU: focused-seeds/AGP/sources/stacks-snapshot-20261011/tags.txt, source lines 9879--9879. Navigation pointer to a statement the counted proof refers to but does not contain; nothing load-bearing is delegated to it.\par\medskip
\begin{lemma}
\label{lemma-finite-module-over-fibre-product}
Let $A, A', B, B', I, M, M', N, \varphi$ be as in
Lemma \ref{lemma-module-over-fibre-product}.
If $N$ finite over $B$ and $M'$ finite over $A'$, then
$N' = N \times_{\varphi, M} M'$ is finite over $B'$.
\end{lemma}

\begin{proof}
We will use the results of
Lemma \ref{lemma-module-over-fibre-product}
without further mention. Choose generators $y_1, \ldots, y_r$ of $N$ over $B$
and generators $x_1, \ldots, x_s$ of $M'$ over $A'$. Using that
$N = N' \otimes_{B'} B$ and $B' \to B$ is surjective we can find
$u_1, \ldots, u_r \in N'$ mapping to $y_1, \ldots, y_r$ in $N$.
Using that $M' = N' \otimes_{B'} A'$ we can find $v_1, \ldots, v_t \in N'$
such that $x_i = \sum v_j \otimes a'_{ij}$ for some $a'_{ij} \in A'$.
In particular we see that the images $\overline{v}_j \in M'$
of the $v_j$ generate $M'$ over $A'$.
We claim that $u_1, \ldots, u_r, v_1, \ldots, v_t$
generate $N'$ as a $B'$-module. Namely, pick $\xi \in N'$. We first choose
$b'_1, \ldots, b'_r \in B'$ such that $\xi$ and $\sum b'_i u_i$ map
to the same element of $N$. This is possible because $B' \to B$
is surjective and $y_1, \ldots, y_r$ generate $N$ over $B$.
The difference $\xi - \sum b'_i u_i$ is of the form $(0, \theta)$
for some $\theta$ in $IM'$. Say $\theta$ is $\sum t_j\overline{v}_j$
with $t_j \in I$. As $J = \Ker(B' \to B)$ maps isomorphically to $I$
we can choose $s_j \in J \subset B'$ mapping to $t_j$.
Because $N' = N \times_{\varphi, M} M'$ it follows
that $\xi = \sum b'_i u_i + \sum s_j v_j$ as desired.
\end{proof}

\end{document}
